% !PRINTED 430
\weberpage{430}

It still remains, when $p$ is a prime and $pp_1$ is a power of $p$
dividing $n-1$, to determine $\varphi(pp_1)$.

If $\alpha$ is a non-primitive congruence root of degree $p^{\pi}$ and
$p^{\pi-1} = p_1$, then, if $\alpha^h$ is the lowest power of $\alpha$
congruent to 1 with respect to the modulus $n$, $h$ is a divisor of $pp_1$,
i.\,e. a power of $p$; but since $h$ is to be smaller than $pp_1$, it
is also a divisor of $p_1$, and consequently $\alpha^{p_1} \equiv 1
\pmod{n}$; thus: a non-primitive congruence root of degree $pp_1$
is at the same time a congruence root of degree $p_1$. Conversely, every
congruence root of degree $p_1$ is at the same time an imprimitive
congruence root of degree $pp_1$. Now since the two congruences
$x^{pp_1} \equiv 1$, $x^{p_1} \equiv 1$ have as many roots as their
degree amounts to, there remain $pp_1 - p_1$ \gesperrt{primitive congruence roots}
of the first, and thus we have
\weq{13}{\varphi(pp_1) = p_1(p-1).}

The function $\varphi(a)$ therefore has the same meaning as in \S.\,132,
and with this the number of primitive roots of the prime
$n$ is at the same time found to be equal to $\varphi(n-1)$. We therefore state the further theorem:

4.\ \gesperrt{A prime $n$ has $\varphi(n-1)$ primitive roots.}

With this the existence of primitive roots for every prime is
proved, and at the same time their number is determined. For finding the
primitive roots we have, to be sure, no other general method
than trial, which can be somewhat shortened by a few
artifices.

If $g$ is a primitive root of the prime $n$, then the residues of the
powers
\weq{14}{1, g, g^2, g^3 \ldots g^{n-2}}
are all different from one another, since, if $g^m \equiv g^{m+\mu}$ held,
$g^{\mu} \equiv 1$ would have to hold, which is not possible so long as $\mu$
is smaller than $n-1$. Hence among the residues of the sequence (14),
among which the residue 0 does not occur, each of the numbers
\weq{15}{1, 2, 3 \ldots n-1}
must be contained, and each only once, or in other words:

5.\ \gesperrt{If $a$ is a number not divisible by $n$, there is
one and only one number $\alpha$ from the sequence of the numbers
$0, 1, 2 \ldots n-2$, by which the congruence}
\weq{16}{g^{\alpha} \equiv a \pmod{n}}
%CONT>>
